\chapter{CANalyzer Integration and Operator GUI}
\label{ch:gui}

Once the firmware and the DBC agreed, the remaining problem was human.
Composing hexadecimal frames in PCAN-View is the right way to validate a
protocol and the wrong way to run a characterisation campaign. This
chapter describes the operator surface that was built on top of Vector
CANalyzer: a CAPL layer that owns the real-time bus, a COM API that
exports system variables, and a Python GUI that owns widgets, gauges and
plots.

\section{Partition of responsibilities}
\label{sec:partition}

Python is a poor real-time CAN stack. Vector CANalyzer is a poor widget
toolkit. The architecture of \cref{fig:gui-arch} takes that observation
literally.

\begin{figure}[ht]
\centering
\begin{tikzpicture}[
  font=\sffamily\footnotesize,
  box/.style={draw=NavyInk, rounded corners=1.5pt, align=center,
              inner sep=6pt, minimum height=12mm},
  arr/.style={-{Latex}, thick, NavyInk}
]
  \node[box, fill=Paper, minimum width=24mm] (mcu) {Stellar-E\\FDCAN};
  \node[box, fill=Paper, minimum width=28mm, right=12mm of mcu]
    (capl) {CANalyzer\\DBC + CAPL};
  \node[box, fill=Paper, minimum width=28mm, right=12mm of capl]
    (com) {COM API\\sysvars};
  \node[box, fill=Paper, minimum width=24mm, right=12mm of com]
    (py) {Python GUI\\widgets};
  \draw[arr] (mcu) -- node[above, font=\scriptsize] {CAN} (capl);
  \draw[arr] (capl) -- (com);
  \draw[arr] (com) -- (py);
  \draw[arr, dashed] ([yshift=-6mm]py.south) -- ++(0,-5mm) -|
    node[near start, below, font=\scriptsize] {write sysvar}
    ([yshift=-6mm]capl.south);
\end{tikzpicture}
\caption{Operator-path architecture. CANalyzer owns timing, DBC encoding
and CAPL; Python owns the graphical interface and talks only to system
variables via \texttt{win32com.client}.}
\label{fig:gui-arch}
\end{figure}

On startup the Python process instantiates
\texttt{Vector.CANalyzer.Application} through \texttt{win32com.client}
\cite{win32com} and obtains the \texttt{App.SystemVariables} namespace.
From that moment on, two directions of traffic exist.

\paragraph{Bus $\rightarrow$ GUI.}
Firmware emits the ten diagnostic frames and, when asked, a
\texttt{REG\_VALUE}. CAPL \texttt{on message} handlers decode them with
the DBC and write system variables. Python polls those variables on a
timer and refreshes gauges, numeric fields and plots.

\paragraph{GUI $\rightarrow$ bus.}
A button, a slider or a ``Set'' action writes one or more system
variables. A CAPL \texttt{on sysvar} handler wakes, checks the ready
flags described below, fills a typed message and calls
\texttt{output()}. CANalyzer encodes the frame from the DBC and
transmits it.

Python never calls a CAN driver. That single restriction is what makes
the GUI interchangeable: a future operator surface in C\# or in CAPL
panels themselves would reuse the same system variables and the same
DBC.

\begin{figure}[ht]
\centering
\includegraphics[width=0.95\textwidth]{canalyzer.png}
\caption{Vector CANalyzer during a session, with the system-variable
namespace that the Python GUI reads and writes through the COM API.}
\label{fig:canalyzer}
\end{figure}

\section{Ready flags for multi-parameter commands}
\label{sec:ready}

A speed ramp has two parameters. If the GUI wrote \texttt{SET\_RAMP = 1}
the moment the operator touched either field, the frame on the bus would
carry a stale duration or a stale final speed. The CAPL layer therefore
keeps three variables per multi-parameter command:

\begin{itemize}[leftmargin=1.6em]
  \item a physical value (\texttt{g\_finalSpeed\_value},
        \texttt{g\_durSpeed\_value});
  \item a ready flag per value (\texttt{is\_finalSpeed\_ready},
        \texttt{is\_durSpeed\_ready});
  \item a trigger (\texttt{SET\_SPEED\_RAMP::SET\_RAMP}).
\end{itemize}

The trigger fires the handler; the handler transmits only when every
ready flag is one. The same pattern is used for current ramps and for PI
writes (numerator and divisor). Incomplete edits stay in the GUI and off
the bus.

\begin{lstlisting}[style=caplstyle,caption={CAPL handler for the speed ramp, with ready-flag guard.},label=lst:capl-speed]
on sysvar SET_SPEED_RAMP::SET_RAMP
{
  message SET_RAMP msgSpeedRamp;

  if (@this == 1 &&
      is_finalSpeed_ready == 1 &&
      is_durSpeed_ready   == 1)
  {
    msgSpeedRamp.FinalSpeed_execrmp = g_finalSpeed_value;
    msgSpeedRamp.Duration_execrmp   = g_durSpeed_value;
    output(msgSpeedRamp);
  }
}
\end{lstlisting}

PI writes are factored into a helper, because three loops times two gains
would otherwise copy the same three assignments nine times:

\begin{lstlisting}[style=caplstyle,caption={CAPL helper that encodes a packed PI SET\_REG.},label=lst:capl-pi]
void send_pi_message(word reg_id, int val, int denom)
{
  message SET_REG msg;
  msg.REGISTER_ID           = reg_id;
  msg.REGISTER_VALUE        = val;
  msg.REGISTER_VALUE_DENOM  = denom;
  output(msg);
}
\end{lstlisting}

\section{Telemetry handlers}
\label{sec:tel-capl}

On the inbound side, CAPL is a set of trivial assignments. The DBC has
already unpacked the frame; the handler's only job is to copy the
physical value into a system variable that Python will poll.

\begin{lstlisting}[style=caplstyle,caption={Inbound CAPL: diagnostics and REG\_VALUE.},label=lst:capl-diag]
on message DIAG_TORQUE_REF
{
  @DIAG_MESSAGES::TORQUE_REF = this.DWord(0);
}

on message REG_VALUE
{
  @GET_REG::K_num = this.GET_K_NUMERATOR;
  @GET_REG::K_den = this.GET_K_DENOMINATOR;
}
\end{lstlisting}

Using \texttt{this.DWord(0)} for the raw diagnostic payload is
intentional. Factors and units are applied either in the DBC signal
definition (when the GUI wants physical units) or in Python (when a plot
needs a custom scale). They are never applied in both places.

\section{Operator panels}
\label{sec:panels}

The GUI is organised as a connection banner plus four tabs. The banner is
always visible: a CAN-bus indicator, a machine-state field, heatsink
temperature, DC-link voltage, and the start / stop / reset-fault
controls. Putting those controls outside the tabs is a safety choice. An
operator who is looking at a PI panel can still stop the motor without
changing tab.

\subsection{Motor control panel}
\Cref{fig:gui-overview} is the first surface the operator sees. Start and
stop map onto the \texttt{EXECUTE\_COMMAND} flags of
\cref{tab:exec-flags}. Temperature and DC-link voltage are bound to
\texttt{DIAG\_TEMPERATURE} and \texttt{DIAG\_BUS\_VOLTAGE}. Machine state
is bound to \texttt{DIAG\_STATE\_MC}; the reset-fault button sets
\texttt{FAULT\_ACK}.

\begin{figure}[ht]
\centering
\includegraphics[width=0.95\textwidth]{gui_overview.png}
\caption{Motor control panel. Connection status, start/stop, thermal and
DC-link readouts, machine state and fault reset live in the top banner
so they remain reachable from every tab.}
\label{fig:gui-overview}
\end{figure}

\subsection{Speed control}
The speed tab (\cref{fig:gui-speed}) exposes the two parameters of
\texttt{SET\_RAMP}: final speed in RPM and ramp duration in milliseconds,
plus an apply button that is the trigger of \cref{lst:capl-speed}. A
gauge shows the reference (\texttt{DIAG\_SPEED\_REF}) against the
measured speed (\texttt{DIAG\_SPEED\_MEAS}) so that a ramp in progress is
visible without reading a plot.

\begin{figure}[ht]
\centering
\includegraphics[width=0.95\textwidth]{gui_speed.png}
\caption{Speed-control tab. The operator sets a final speed and a
duration; the gauge compares reference and actual RPM in real time.}
\label{fig:gui-speed}
\end{figure}

\subsection{PI controller}
The PI tab (\cref{fig:gui-pi}) is a direct view of
\cref{sec:pi-pack}. Three groups---speed, flux, torque---each offer $K_p$
and $K_i$ as a numerator and a power-of-two divisor. ``Set'' calls
\texttt{send\_pi\_message} with the corresponding register identifier.
``Read'' issues \texttt{GET\_REG} and fills the same fields from
\texttt{REG\_VALUE}, which is how an operator confirms that a write
round-tripped.

\begin{figure}[ht]
\centering
\includegraphics[width=0.95\textwidth]{gui_pi.png}
\caption{PI-controller tab. Each of the three loops exposes $K_p$ and
$K_i$ as a numerator / $2^n$ pair, matching the packed 32-bit register
layout.}
\label{fig:gui-pi}
\end{figure}

\subsection{Current control}
The current tab (\cref{fig:gui-current}) is the operator counterpart of
\cref{sec:id-iq,sec:iq-ramp,sec:id-ramp}. $I_q$ and $I_d$ can be written
as instantaneous set-points or as ramps with their own durations. Two
strip charts show reference against measurement for each axis, which is
the quantity a current-loop tuning session actually looks at.

\begin{figure}[ht]
\centering
\includegraphics[width=0.95\textwidth]{gui_current.png}
\caption{Current-control tab. Independent $I_q$ and $I_d$ set-points and
ramps, with live plots of reference and measured current on each axis.}
\label{fig:gui-current}
\end{figure}

\subsection{Live charts}
The last tab (\cref{fig:gui-charts}) repeats the speed gauge and the two
current plots at a larger size, without the entry fields. It is the
surface used once a test is running and the operator's job is to watch,
not to command.

\begin{figure}[ht]
\centering
\includegraphics[width=0.95\textwidth]{gui_charts.png}
\caption{Live-charts tab: speed overview together with the $I_q$
(torque) and $I_d$ (flux) current plots.}
\label{fig:gui-charts}
\end{figure}

\section{What the GUI does not do}
A short negative list, because the temptation to grow an operator surface
into a second dyno controller is strong and was resisted.

\begin{itemize}[leftmargin=1.6em]
  \item The GUI does not implement a safety PLC. Emergency stop remains a
        hardware input to the inverter and to the dyno.
  \item The GUI does not log to disk. CANalyzer already records traces in
        a format that the rest of the laboratory toolchain consumes.
  \item The GUI does not re-implement DBC packing. If a signal scale is
        wrong, the fix is in the DBC, not in a Python constant.
  \item The GUI does not talk to the MCU except through CANalyzer. There
        is no second CAN adapter in the Python process, so there is no
        second source of frames to debug.
\end{itemize}

Those restrictions are what keep the three-artefact rule of
\cref{ch:intro} intact: firmware, DBC, operator surface, in that order of
authority.
